\Ch{Terms and definitions}{Terms}

\p This document uses terms consistent with their definitions in \gls{isoC} and
\gls{isoCPP}. In cases where the definitions are unclear, or where this document
diverges from \gls{isoC} and \gls{isoCPP}, the definitions in this clause, and
the attached glossary (\ref{main}) supersede other sources.

\Sec{Common Definitions}{Terms.Common}

\p The following definitions are consistent between \acrshort{hlsl} and the
\gls{isoC} and \gls{isoCPP} specifications, however they are included here for
reader convenience.

\Sub{Correct Data}{Terms.CorrectData}
\p Data is correct if it represents values that have specified or unspecified
but not undefined behavior for all the operations in which it is used. Data that
is the result of undefined behavior is not correct, and may be treated as
undefined.

\Sub{Diagnostic Message}{Terms.Diags}
\p An implementation defined message belonging to a subset of the
implementation's output messages which communicates diagnostic information to
the user.

\Sub{Ill-formed Program}{Terms.IllFormed}
\p A program that is not well-formed, for which the implementation is expected
to return unsuccessfully and produce one or more diagnostic messages.

\Sub{Implementation-defined Behavior}{Terms.ImpDef}
\p Behavior of a well-formed program and correct data which may vary by the
implementation, and the implementation is expected to document the behavior.

\Sub{Implementation Limits}{Terms.ImpLimits}
\p Restrictions imposed upon programs by the implementation of either the
compiler or runtime environment. The compiler may seek to surface
runtime-imposed limits to the user for improved user experience.

\Sub{Undefined Behavior}{Terms.Undefined}
\p Behavior of invalid program constructs or incorrect data for which this
standard imposes no requirements, or does not sufficiently detail.

\Sub{Unspecified Behavior}{Terms.Unspecified}
\p Behavior of a well-formed program and correct data which may vary by the
implementation, and the implementation is not expected to document the behavior.

\Sub{Well-formed Program}{Terms.WellFormed}
\p An \acrshort{hlsl} program constructed according to the syntax rules,
diagnosable semantic rules, and the One Definition Rule.

\Sec{General Terms}{Terms.General}

\p The following terms are specific to \acrshort{hlsl} and are used throughout
the specification.

\Sub{Runtime Implementation}{Terms.Runtime}

\p A runtime implementation refers to a full-stack implementation of a software
runtime that can facilitate the execution of \acrshort{hlsl} programs. This
broad definition includes libraries and device driver implementations. The
\acrshort{hlsl} specification does not distinguish between the user-facing
programming interfaces and the vendor-specific backing implementation.

\Sub{Host and Device}{Terms.HostDevice}

\p \acrshort{hlsl} is a data-parallel programming language designed for
programming auxiliary processors in a larger system. In this context the
\textit{host} refers to the primary processing unit that runs the application
which in turn uses a runtime to execute \acrshort{hlsl} programs on a supported
\textit{device}. There is no strict requirement that the host and device be
different physical hardware, although they commonly are. The separation of host
and device in this specification is useful for defining the execution and memory
model as well as specific semantics of language constructs.

\Sec{\acrshort{spmd} Terminology}{Terms.SPMD}

\p \acrshort{hlsl} is a \acrfull{spmd} programming language. The following terms
are used to describe the execution model of \acrshort{spmd} programs and the
semantics of language constructs that are specific to \acrshort{spmd}
programming.

\Sub{\gls{lane}}{Terms.Lane}

\p A \gls{lane} represents a single computed element in an \acrshort{spmd}
program. In a traditional programming model it would be analogous to a thread of
execution, however it differs in one key way. In multi-threaded programming,
threads advance independent of each other. In \acrshort{spmd} programs, a group
of \gls{lane}s may execute instructions in lockstep because each instruction may
be a \acrshort{simd} instruction computing the results for multiple \gls{lane}s
simultaneously, or synchronizing execution across multiple \gls{lane}s or
\gls{wave}s. A \gls{lane} has an associated \textit{lane state} which denotes
the execution status of the lane (\ref{Terms.LaneState}).

\Sub{\gls{wave}}{Terms.Wave}

\p A grouping of \gls{lane}s for execution is called a \gls{wave}. The size of a
\gls{wave} is defined as the maximum number of \textit{active} \gls{lane}s the
\gls{wave} supports. \gls{wave} sizes vary by hardware architecture, and are required
to be powers of two. The number of \textit{active} \gls{lane}s in a \gls{wave}
can be any value between one and the \gls{wave} size.

\p Some hardware implementations support multiple \gls{wave} sizes. There is no
overall minimum wave size requirement, although some language features do have
minimum \gls{lane} size requirements.

\p \acrshort{hlsl} is explicitly designed to run on hardware with arbitrary
\gls{wave} sizes. Hardware architectures may implement \gls{wave}s as
\acrfull{simt} where each thread executes instructions in lockstep. This is not
a requirement of the model. Some constructs in \acrshort{hlsl} require
synchronized execution. Such constructs will explicitly specify that
requirement.

\Sub{\gls{quad}}{Terms.Quad}

\p A \gls{quad} is a subdivision of four \gls{lane}s in a \gls{wave} which are
computing adjacent values. In pixel shaders a \gls{quad} may represent four
adjacent pixels and \gls{quad} operations allow passing data between adjacent
\gls{lane}s. In compute shaders quads may be one or two dimensional depending
on the workload dimensionality. Quad operations require four active \gls{lane}s.

\Sub{\gls{threadgroup}}{Terms.Group}

\p A grouping of \gls{lane}s executing the same shader to produce a combined
result is called a \gls{threadgroup}. \gls{threadgroup}s are independent of
\acrshort{simd} hardware specifications. The dimensions of a \gls{threadgroup}
are defined in three dimensions. The maximum extent along each dimension of a
\gls{threadgroup}, and the total size of a \gls{threadgroup} are implementation
limits defined by the runtime and enforced by the compiler. If a
\gls{threadgroup}'s size is not a whole multiple of the hardware \gls{wave}
size, the unused hardware \gls{lane}s are implicitly inactive.

\p If a \gls{threadgroup} size is smaller than the \gls{wave} size , or if the
\gls{threadgroup} size is not an even multiple of the \gls{wave} size, the
remaining \gls{lane} are \textit{inactive} \gls{lane}s.

\Sub{\gls{dispatch}}{Terms.Dispatch}

\p A grouping of \gls{threadgroup}s which represents the full execution of a
\acrshort{hlsl} program and results in a completed result for all input data
elements.

\Sub{\gls{lane} States}{Terms.LaneState}

\p \gls{lane}s may be in four primary states: \textit{active}, \textit{helper},
\textit{inactive}, and \textit{predicated off}.

\p An \textit{active} \gls{lane} is enabled to perform computations and produce
output results based on the initial launch conditions and program control flow.

\p A \textit{helper} \gls{lane} is a lane which would not be executed by the
initial launch conditions except that its computations are required for adjacent
pixel operations in pixel fragment shaders. A \textit{helper} \gls{lane} will
execute all computations but will not perform writes to buffers, and any outputs
it produces are discarded. \textit{Helper} lanes may be required for
\gls{lane}-cooperative operations to execute correctly.

\p An \textit{inactive} \gls{lane} is a lane that is not executed by the initial
launch conditions. This can occur if there are insufficient inputs to fill all
\gls{lane}s in the \gls{wave}, or to reduce per-thread memory requirements or
register pressure.

\p A \textit{predicated off} \gls{lane} is a lane that is not being executed due
to program control flow. A \gls{lane} may be \textit{predicated off} when
control flow for the \gls{lane}s in a \gls{wave} diverge and one or more lanes
are temporarily not executing.

\p The diagram blow illustrates the state transitions between \gls{lane} states:

\begin{center}
\begin{figure}
  \centering
  \caption{State transitions for \gls{lane} states}
  \begin{tikzpicture}[shorten >=1pt,node distance=3cm,auto, squarednode/.style={rectangle,minimum size=7mm}]
    \node[squarednode,state,initial above] (active) {$active$};
    \node[squarednode,state,initial above] (inactive)[right of=active] {$inactive$};
    \node[squarednode,state] (helper)[below of=active] {$helper$};
    \node[squarednode,state] (off1)[below left of=active] {off};
    \node[squarednode,state] (off2)[right of=helper] {off};


    \path[->] (active) edge node {discard} (inactive);
    \path[->] (active) edge node {discard} (helper);
    \path[->] (active) edge node {branch} (off1);
    \path[->] (off1) edge node {} (active);

    \path[->] (helper) edge node {branch} (off2);
    \path[->] (off2) edge node {} (helper);
  \end{tikzpicture}
\end{figure}
\end{center}

\Sub{Uniformity}{Terms.Uniformity}

\p Uniformity is a property of either a value or a control flow construct.
Uniformity can be discussed in terms of any grouping of \gls{lane}s (e.g.
\gls{quad}, \gls{threadgroup}, \gls{wave}, \gls{dispatch}), if a grouping of
\gls{lane}s is not explicitly mentioned, it is assumed to be a \gls{wave}.

\p When referring to a value, uniformity means that the value is the same across
all \gls{lane}s in the grouping. When referring to a control flow construct,
uniformity means that all \gls{lane}s in the grouping are executing the same
path through the program's control flow.

\p Uniformity can be a static property that is determined at compile time, or a
dynamic property that is determined at runtime.

\Sub{Divergence}{Terms.Divergence}

\p Divergence is a dynamic property of program execution where one or more
\gls{lane}s in a \gls{wave} are executing different paths through the program's
control flow. Control flow that causes divergence is said to be
\textit{divergent}.

\Sub{Convergence}{Terms.Convergence}

\p Convergence is a property of an operation that requires the \gls{lane}s
executing the operation to synchronize in order to produce correct results.
Operations requiring convergence are said to be \textit{convergent}.

\Sec{Memory Spaces}{Terms.MemorySpaces}

\p Memory spaces refer to the different types of memory that programs can
access. Each memory space has different properties and semantics for how it can
be accessed and used.
